\subsection{What Is Included in a Policy Bill?}
\label{sec:resources}
Checking is called \emph{calibration} in the execution records; the two terms denote the same phase. A request, or model call, processes one prompt. A standalone policy pays for CPU grouping, its own singleton checking references, all issued probe chunks, and complete execution at the selected size or singleton fallback. S1/S2 and independent validation score the policy for the study; they are charged to the research ledger, rather than to that standalone policy. Each policy starts with an empty prefix cache and may reuse only its own checking-to-execution prefix computation. Cache OFF is a separately executed configuration; operator and response memoization are disabled.

For requests, the accounting is simply checking plus execution: Sports uses $160+256=416$, compared with 2,048 for S1. Token accounting needs another distinction. Let $P$ be all input tokens written in requests, $O$ all generated output tokens, and $K$ the input tokens the server reports as served from its prefix cache. The totals cover the same policy phases. We retain both
\begin{equation}
\underbrace{P+O}_{\text{raw input-plus-output tokens}},\qquad
\underbrace{U=P-K}_{\text{uncached input tokens}}.
\label{eq:costs}
\end{equation}
Raw tokens count the prompt and response volume. Uncached input subtracts reused input computation and is a proxy for fresh prefill work; output generation remains separate~\cite{apc}. Neither quantity is a dollar bill or a runtime prediction. Missing cached-token counters remain unknown with bounds, rather than being replaced by zero. Sports will illustrate the complete raw-token bill; matched cache ON/OFF cases will illustrate the second cost measure.

Our primary \emph{joint benefit criterion} asks for all four conditions: (i) a complete native result; (ii) actual batching, $k>1$; (iii) at most 5\% row disagreement with each of S1/S2; and (iv) fewer requests \emph{and} fewer raw input-plus-output tokens than S1. A fallback can preserve reference labels, but it is not a successful batch deployment. Source accuracy and uncached input are reported separately. The strict-vector sensitivity evaluates the same result and bill using complete strict labels; it keeps unavailable cases in the 12-case denominator.

Elapsed time starts after a verified cache reset and includes CPU grouping, client isolation, durable tracing, transport, checking and execution. Input preparation, model loading and the reset are outside the timer. Shared-server times are descriptive; the plan comparisons below use explicit requests and token counters.
